% 四值不等式的逐格计数：传感器实例取 X={a,b}、Y={b,c}。
% (a) 每格计入“X覆盖它”与“Y覆盖它”的次数，合计 F(X)+F(Y)=12+16=28。
% (b) 每格计入“X∪Y覆盖它”与“X∩Y={b}覆盖它”的次数，合计 F(X∪Y)+F(X∩Y)=16+8=24。
% 第1--4列由公共部分 b 覆盖（蓝底），两侧都计2；第5--6列由 X 中的 a 与 Y 中的 c 分别覆盖、b 未覆盖（黄底），
% 左侧计2、右侧计1，共4格，正是差值 28-24=4；第7--8列只由 Y 覆盖，两侧都计1。数据由 scratchpad 中 verify.py 复算。
\begin{tikzpicture}[x=6.2mm,y=6.2mm,line width=.6pt,font=\figurefont\fontsize{8}{10}\selectfont]
  \foreach \yy/\n/\m in {2.7/2/1,0/1/1}{
    \begin{scope}[shift={(0,\yy)}]
      \foreach \c in {1,...,4}{\foreach \y in {0,1}{
        \fill[FigureBlueFill] (\c-1+.04,\y+.04) rectangle (\c-.04,\y+.96);}}
      \foreach \c in {5,6}{\foreach \y in {0,1}{
        \fill[FigureYellowFill] (\c-1+.04,\y+.04) rectangle (\c-.04,\y+.96);}}
      \foreach \c in {1,...,8}{\foreach \y in {0,1}{
        \draw[FigureGrid,line width=.5pt] (\c-1,\y) rectangle (\c,\y+1);}}
    \end{scope}
  }
  % (a) 分别评价
  \begin{scope}[shift={(0,2.7)}]
    \foreach \c/\v in {1/2,2/2,3/2,4/2,5/2,6/2,7/1,8/1}{\foreach \y in {0,1}{
      \node[FigureInk] at (\c-.5,\y+.5) {\v};}}
    \node[anchor=west,FigureInk] at (8.3,1) {$F(X)+F(Y)=28$};
    \node[anchor=east,FigureInk] at (-0.3,1) {(a) 分别评价};
  \end{scope}
  % (b) 并与交
  \foreach \c/\v in {1/2,2/2,3/2,4/2,5/1,6/1,7/1,8/1}{\foreach \y in {0,1}{
    \node[FigureInk] at (\c-.5,\y+.5) {\v};}}
  \node[anchor=west,FigureInk] at (8.3,1) {$F(X\cup Y)+F(X\cap Y)=24$};
  \node[anchor=east,FigureInk] at (-0.3,1) {(b) 并与交};
  \foreach \c in {1,...,8}{\node[FigureMuted,font=\FigureNoteFont\fontsize{7.5}{9}\selectfont] at (\c-.5,-0.5) {\c};}
  \node[FigureMuted,anchor=east,font=\FigureNoteFont\fontsize{7.5}{9}\selectfont] at (-0.3,-0.5) {列号};
\end{tikzpicture}
